\section{Methods}\label{sec:methods}

We conducted a systematic review of robotic active matter following the
PRISMA 2020 reporting framework \todocite{Page et al., PRISMA 2020
statement, BMJ 2021}. The review was designed as a living evidence
synthesis, with periodic literature searches and versioned updates
\todocite{Elliott et al., living systematic review methodology, 2017}.
Our central aim was to distinguish experimental demonstrations of
collective behavior from quantitative tests of the physical mechanisms
underlying that behavior. Hence, we treated robotic active matter as
an experimental research field rather than as a subset of swarm-robotics
engineering.

In the following, we describe the eligibility criteria, search and
selection procedures, extraction and adjudication methods, assessment
of study quality, and synthesis strategy. We also report the validation
of the language-model-assisted screening procedure. The numerical
results refer to the evidence snapshot of October 17, 2026.

\subsection{Eligibility criteria}\label{sec:methods-eligibility}

We defined the population as physical collectives containing at least
two interacting artificial agents that were self-propelled or actively
driven. Eligible platforms included macroscopic swarm robots,
vibration-driven robots, modular systems, robotic micro- and
nano-agents, actively controlled synthetic colloids, and robotic
metamaterials. Two-agent experiments were eligible because pairwise
interactions can reveal the elementary mechanisms from which collective
behavior emerges.

The decisive inclusion criterion was a quantitative observable
describing collective behavior or interaction. Examples include
polarization, spatial correlations, cluster statistics, density
fluctuations, transport coefficients, phase diagrams, and mechanical
response. We accepted primary experimental research published as
journal articles, conference papers with results, or preprints.
Theoretical and simulation studies qualified only when they included
quantitative comparison with robotic hardware experiments.
Publication date was unrestricted, and we considered English-language
reports.

This definition deliberately separates robotic active matter from
general swarm robotics. A system that achieves coverage, navigation,
or formation control is not necessarily an experimental realization
of a collective physical phenomenon. We therefore excluded studies
reporting only engineering performance, such as task completion time
or algorithmic success rates, without a physical collective
observable. We also excluded single-agent investigations,
simulation-only studies without experimental comparison, studies
involving exclusively biological agents, and non-primary
publications.

Synthetic active colloids required an additional distinction.
Autonomous chemical propulsion or uniform-field actuation alone did
not establish eligibility. Colloidal studies qualified when
programmable spatial or temporal actuation, feedback, or learning
controlled the system. This restriction preserved a meaningful
connection between robotic programmability and active-matter physics.

We used six mutually exclusive exclusion codes. When several
criteria applied, we assigned the first applicable code in the
prespecified precedence order: non-primary research (E5),
biological-only systems (E4), uncontrolled synthetic colloids (E6),
single-agent systems (E1), theory or simulation without qualifying
hardware evidence (E2), and engineering-only swarm robotics (E3).
This hierarchy ensured that each excluded report received exactly
one primary reason.

\subsection{Information sources and search strategy}
\label{sec:methods-search}

We searched OpenAlex, Crossref, arXiv, the Europe PMC index of
bioRxiv preprints, PubMed, IEEE Xplore, HAL, and NASA ADS.
These sources were chosen to cover the intersection of statistical
physics, soft condensed matter, robotics, control, and engineering.
SemOpenAlex provided a supplementary metadata and citation-lookup
interface, rather than an independent search source.
Google Scholar, citation chasing, and journal hand-searching were
specified as supplementary discovery procedures.

The search covered database inception through the snapshot date.
Its conceptual structure combined three routes into the literature:
active matter coupled to robotic agents, controlled synthetic colloids,
and established composite expressions for relevant systems.
Schematically, the Boolean logic was

\begin{quote}
\raggedright\small
(active matter and robotic agent)
or (colloid and control)
or composite phrases
\end{quote}

These expressions denote concept blocks rather than the complete
database-specific executable queries. Each source required its own
syntax. Crossref and the bioRxiv index, for example, required
separate phrase searches followed by local filtering.
The protocol specifies that complete search strings, execution
timestamps, record counts, raw-response checksums, and filter versions
are retained in dated strategy files.

A pilot sensitivity check on October 11, 2026 retrieved all
12 prespecified seed papers through at least one primary source.
The search design also specified peer review using the PRESS 2015
checklist \todocite{PRESS 2015 guideline for peer reviewing electronic
search strategies}. The supplied records establish the intended
procedure but do not independently document completion of PRESS
review or provide all executable source-specific search strings.

We deduplicated records by canonicalized DOI and then by normalized
title similarity, using a Jaccard threshold of 0.9.
Merged records retained their source identifiers. Records were
subsequently screened at title/abstract level and, where possible,
assessed using retrieved full texts.

\subsection{Study selection and human validation}
\label{sec:methods-screening}

Screening used three large language models accessed through the
NHR@FAU language-model gateway: GPT-OSS-120B (R1),
Magistral-Small-2509 (R2), and DeepSeek-V4-Flash (R3).
Each reviewer received the same screening instructions in an
independent context, without access to the other reviewers'
decisions. The nominal sampling temperature was zero.
Each assessment recorded an inclusion decision, exclusion reason,
confidence, and a short rationale, together with model identity,
prompt checksum, and timestamp.

The decision rule was a majority of three reviewers. To reduce
computational cost, R1 and R2 evaluated every record, while R3
was called whenever its vote could change inclusion.
Additional R3 assessments resolved relevant full-text
exclusion-code disagreements and supported a deterministic audit
sample. Thus, the resulting inclusion decision was intended to
equal that of a complete three-reviewer majority, although not
every record received three votes.

A record without an abstract proceeded toward full-text assessment
unless all three reviewers excluded it based on its title.
Unretrievable full texts were handled under a protocol amendment:
eligibility could be assessed from the abstract when no report
could be obtained. Such studies were explicitly flagged and
retained for a separate sensitivity analysis. Conference abstracts
without quantitative results remained ineligible.

Note that agreement among language models is not equivalent to
validation against an independent reference. We therefore evaluated
screening decisions against an author-labelled random sample of
200 title/abstract records. The author initially labelled the
sample without seeing panel decisions and subsequently revised
provisional labels after inspecting the records. Both the first
blind labels and the final labels were retained.

Against the author's final labels, the panel correctly identified
20 of 20 eligible records and 177 of 180 ineligible records.
Sensitivity was 1.000, specificity was 0.983, and Cohen's
$\kappa$ was 0.922. There were three false inclusions and no
false exclusions. The prespecified sensitivity threshold of
0.95 was therefore met under the final-label evaluation.

Individual models showed substantial differences.
R1 achieved sensitivity 0.950, specificity 0.950, and
$\kappa=0.764$. R2 achieved sensitivity 0.600,
specificity 0.994, and $\kappa=0.704$.
We additionally evaluated Claude Opus as an independent blind
rater on the same sample. It achieved sensitivity 0.550,
specificity 1.000, and $\kappa=0.687$ against the final
author labels.

The original blind author labels provide an important
qualification. Against these first labels, panel sensitivity
was 0.290, specificity 0.917, and $\kappa=0.232$.
The corresponding values were 0.323, 0.893, and 0.225
for R1; 0.161, 0.953, and 0.149 for R2; and
0.161, 0.964, and 0.171 for Claude Opus.

In particular, the original blind author labels classified
31 records as eligible, whereas the final labels classified
20 as eligible. The first-label analysis is therefore not an
independent confirmation of the final-label results.
The improvement in measured panel sensitivity depends on the
author's subsequent corrections. We report both evaluations
rather than treating the final-label comparison as evidence of
unconditional screening reliability.

\begin{figure}[t]
\centering
\includegraphics[width=\columnwidth]{figures/validation_agreement}
\caption{Validation of language-model-assisted title/abstract
screening against the author's final labels in a random sample
of 200 records. Sensitivity, specificity, and Cohen's
$\kappa$ are shown for the majority panel, its individual
reviewers R1 and R2, and Claude Opus as an additional
independent blind rater. The dashed line indicates the
prespecified sensitivity threshold of 0.95. The final-label
comparison is supplemented in the text by results against
the author's initial blind labels, which differed
substantially.}
\label{fig:validation}
\end{figure}

Figure~\ref{fig:validation} illustrates the final-label
comparison. The panel's advantage in sensitivity over the
individual models supports majority-based screening under
this reference standard. Yet, the difference between blind
and final author labels shows that validation remains
sensitive to the construction of the reference set.

\subsection{Evidence-based conflict adjudication}
\label{sec:methods-adjudication}

Initial screening decisions were not treated as immutable.
The protocol included author adjudication of conflicting
full-text decisions and unresolved extraction fields.
An amendment dated October 15, 2026 introduced an
evidence-based procedure to make these judgments auditable.

For each disputed decision, the original reviewers were
asked to reconsider their answers and provide up to four
short passages copied from the source report.
Each quotation was checked automatically against the
extracted text and classified according to its degree of
verbatim correspondence. Unsupported quotations could
therefore be distinguished from passages present in the
document.

The revised model answers were displayed as evidence,
not counted as additional votes. This distinction matters:
otherwise repeated consideration by the same model could
create an artificial majority without adding an independent
reviewer. The intended adjudication was based on the
documentary support for competing claims.

On October 17, 2026, the author delegated the remaining
423 full-text split decisions and 167 extraction-field
conflicts to a separate Claude Opus judge. The judge
received the reviewers' claims and checked quotations.
When the evidence was contradictory or insufficient,
it was instructed to inspect the full text directly.
Unlike the original panel, the judge was explicitly
instructed to weigh documentary evidence rather than
reviewer confidence or vote counts.

For the full-text splits, the judge agreed with the
panel majority on 362 decisions and reversed it on 59.
Two cases remained undecided. Agreement was therefore
362 of 421 decided cases, or 86.0\%.
Of the reversals, nine changed exclusion to inclusion
and 50 changed inclusion to exclusion.
The judge consulted full text directly in 133 cases.
Three additional records were flagged because the
stored full text belonged to a different publication;
the supplied adjudication record does not establish
their subsequent resolution.

This pattern reveals an asymmetry worth noting.
The judge more often overturned positive than negative
panel decisions. Yet, this does not independently prove
that the judge's decisions were correct. Both stages
relied on language-model interpretation, and the
judge's output was not a new human gold standard.
It was instead a documented adjudication layer,
subject to later author correction.

For the 167 extraction-field conflicts, the judge
selected R3's value in 116 cases, R1's in 17,
R2's in 17, and an independently chosen value in 17.
Eighteen fields were set to null, and 68 decisions
required direct consultation of the full text.
These totals describe adjudication choices, not
independent measurements of extractor accuracy.

The protocol gives precedence to existing author
decisions. Judge decisions may replace panel decisions
but may not override author labels. All decisions and
rationales are specified for retention in an audit ledger.
The associated judge flags also identify possible
duplicate records and source-document mismatches.
These issues represent limitations of the snapshot
rather than evidence that all bibliographic ambiguities
have been eliminated.

\subsection{Data extraction}\label{sec:methods-extraction}

We extracted study characteristics using a structured
form with controlled vocabularies. The principal
dimensions were platform class, agent number and size,
propulsion mechanism, interaction modality, control
architecture, experimental environment, observed
phenomena, reported observables, and theoretical
frameworks used for comparison.

Platform classes distinguished vibration-driven robots,
wheeled or legged swarms, modular robots, smarticles
and granular robots, micro- and nano-agents,
feedback-controlled colloids, robotic active solids,
and other systems. Interactions could be steric,
hydrodynamic, magnetic, optical or sensing-based,
communicative, programmed, or mediated by the
environment. Multiple interaction categories and
collective phenomena could be assigned to one study.

Control architectures were coded as passive or
uncontrolled rules, on-board rules, external feedback,
learned policies, central control, or other.
Theoretical comparisons were classified by framework
and by whether agreement was quantitative,
qualitative, or absent. We also extracted the number
of independent experimental runs and explicit
availability of data and code.

The extraction procedure specified independent
assessment by R1 and R3. Disputed fields were sent
to R2 and settled by majority where possible.
For multi-label fields, majority voting was applied
to each label separately rather than to the complete
set of labels. Numerical values were compared using
prespecified tolerances: 10\% for counts and 25\%
for agent size. Single-valued conflicts without a
majority were flagged for adjudication.

Missing information was not imputed. In particular,
an unreported experimental repetition count cannot
be interpreted as evidence that only one experiment
was performed. Similarly, code and data availability
were marked positive only when a repository or
archive was explicitly identified.

At this snapshot, 462 included studies had a
full-text assessment basis and 30 were assessed
from abstracts alone. For the latter, extraction
necessarily reflects the more limited information
available in an abstract. We therefore distinguish
missing reporting from missing experimental practice
and assess the influence of abstract-only inclusion
in sensitivity analyses.

\subsection{Risk of bias and study quality}
\label{sec:methods-bias}

No validated risk-of-bias instrument was identified
in the protocol for the heterogeneous experiments
covered by robotic active matter. We therefore
specified a domain-based quality assessment designed
around the principal sources of uncertainty in
collective-behavior experiments.

The six domains were independent replication (D1),
characterization of agent heterogeneity (D2),
finite-size and boundary effects (D3),
measurement validity (D4), data and code
availability (D5), and quantitative comparison
with theory (D6). Each domain received a rating
of low risk, some concerns, or high risk.
Three reviewer assessments were combined by
majority; a three-way disagreement resulted in
the intermediate rating.

These categories require careful interpretation.
A high-risk rating for quantitative theory
comparison does not necessarily imply that the
reported experimental observation is false.
It means that the study provides limited grounds
for testing a proposed theory quantitatively.
Similarly, failure to report independent runs
weakens reproducibility without establishing
that replication did not occur.

For synthesis, we emphasize D1 and D6 because
the central review questions concern reproducible
collective behavior and quantitative confrontation
with theory. We also retain the other domains,
since agent heterogeneity, boundaries, and
measurement procedures may determine whether an
observed pattern reflects a collective mechanism
or an experimental artifact.

\subsection{Effect measures and synthesis}
\label{sec:methods-synthesis}

A conventional meta-analysis was not appropriate.
The included studies report heterogeneous
systems and observables, and their effect measures
are generally not commensurable. Polar order,
cluster size, velocity correlations, transport,
and elastic response cannot be combined into
a single meaningful effect estimate.

Instead, we use a structured narrative synthesis
supported by evidence maps. The primary mappings
connect platform classes to observed phenomena,
and phenomena to theoretical frameworks.
A study may contribute to several categories.
Consequently, category totals need not sum to
the number of unique included studies.

We distinguish observation from explanation.
A measured collective pattern establishes that
a phenomenon occurs under particular experimental
conditions. A model reproducing the pattern
qualitatively provides a possible explanation.
A quantitative comparison across independently
varied parameters provides stronger evidence
for the underlying physical mechanism.

This hierarchy guides the synthesis without
reducing it to a vote count of favorable papers.
We consider the number of independent platforms,
the diversity of experimental groups,
measurement quality, replication, and the
specific observables tested. We also distinguish
agreement achieved by fitting experimental
parameters from predictions tested outside
the calibration regime whenever the extracted
information permits that distinction.

Sensitivity analyses are specified for studies
with high risk in D1 or D6, and separately for
studies evaluated only from abstracts. These
analyses ask whether the qualitative interpretation
of the evidence changes when the least secure
contributions are removed. The protocol specifies
these comparisons; their substantive consequences
are considered in the synthesis and discussion.

\subsection{Reporting bias and certainty of evidence}
\label{sec:methods-certainty}

We assess possible reporting bias by comparing
findings in preprint-only and peer-reviewed
studies. Such a comparison can reveal differences
in the reporting of positive observations,
quantitative validation, or reproducibility
information. Yet, publication status is not
itself a reliable measure of experimental quality,
and differences cannot automatically be
attributed to publication bias.

For each substantive synthesis statement, the
protocol specifies a certainty judgment of
high, moderate, or low. The assessment emphasizes
independent reproduction across platforms and
research groups, together with the replication
and theory-comparison domains D1 and D6.

A large number of demonstrations on closely
related platforms does not automatically
establish high certainty. Conversely, a smaller
set of independently reproduced quantitative
tests may provide stronger support for a general
physical principle. Our judgments therefore
concern the strength and transferability of
specific conclusions, rather than an overall
ranking of publications.

The language-model-assisted design introduces
an additional review-level limitation.
Screening, extraction, and conflict resolution
may share systematic interpretation errors,
even when individual model outputs are independent.
We address this through the validation sample,
documented evidence quotations, author
override provisions, and transparent reporting
of unresolved cases. These safeguards improve
auditability but do not eliminate correlated
model errors.

\subsection{Living-review procedure, registration,
and transparency}\label{sec:methods-living}

The review is designed for monthly search updates.
New records are to undergo the same screening
and extraction procedures, with dated evidence
snapshots and an updated PRISMA flow diagram.
Material changes in the synthesis should trigger
a new manuscript version. Candidate triggers
include the emergence of a new platform class,
a previously unrepresented phenomenon, or a
change in the certainty assigned to a central
conclusion.

The living structure is particularly important
for robotic active matter because programmable
platforms can introduce new physical interactions
and experimental control regimes. Hence, an
update must reconsider theoretical conclusions
rather than merely append newly published studies.

The supplied protocol describes registration
with OSF Registries but retains a draft status.
A completed registration identifier and
confirmation of protocol freezing are not
documented in the supplied materials.
We therefore report the registration as
[OSF registration status and identifier to be
confirmed], rather than implying prospective
registration that cannot be verified.
The dated amendment log records changes to
screening, evidence-based adjudication,
extraction reconciliation, and validation
definitions.

The protocol specifies public availability of
pipeline code, search strategies, screening
decisions, extracted data, and dated snapshots,
while restricting redistribution of
subscription-only full texts. The archival
repository and persistent access information
are [repository URL and persistent identifier
to be confirmed]. No study authors were
contacted for missing information in the
initial review.

Funding and institutional support are
[funding and support information to be supplied].
Competing interests are
[competing-interest declarations to be supplied].
These items cannot be inferred from the
review data.

Finally, we disclose the use of artificial
intelligence in the research and writing
workflow. Language models supported screening,
data extraction, evidence gathering, and
conflict adjudication as described above.
Manuscript prose was drafted using an AI
system, Codex, for subsequent author review
and revision. The protocol specifies that
the drafting exchanges are archived in the
public repository. This disclosure is intended
to support accountability and compliance
with the journal's AI-use requirements
\todocite{current APS policy on generative
artificial intelligence and manuscript disclosure}.

The responsibility for the scientific
interpretation, the final inclusion set,
and the published conclusions remains
with the authors. The review's central
methodological principle is therefore
to preserve the distinction between
a model-generated judgment, the
documentary evidence supporting it,
and a conclusion that has undergone
human scientific assessment.
